\problem{}
\probitem{}

Let an element of the Lorentz group be represented by, \begin{align}U\qty(\Lambda) = \exp\qty(\frac\im 2\omega_{\mu\nu}M^{\mu\nu})\end{align} this is indirectly defined through 
the defining representation of the Lorentz group, \begin{align}\tensor{\Lambda}{^\alpha_\beta}\qty(\omega)%\rightarrow -\im\dv{\tensor{\Lambda}{^\alpha_\beta}}{\omega_{\mu\nu}}\eval_{\omega=0} = -\frac\im 2 \qty(\tensor{g}{^\alpha^\mu}\tensor{g}{_\beta^\nu}-\tensor{g}{^\alpha^\nu}\tensor{g}{_\beta^\mu})
\end{align} 
using our knowledge of the composition law of two transformations,
\begin{align}
    U\qty(\Lambda_1)U\qty(\Lambda_2)U^{-1}\qty(\Lambda_1) &=U(\Lambda_1\Lambda_2 \Lambda_1^{\textnormal T})\\
    U\qty(\Lambda_1)U\qty(\mathbbm 1+\omega_2)U^{-1}\qty(\Lambda_1) &=U(\Lambda_1 \qty(\mathbbm 1+\omega_2) \Lambda_1^{\textnormal T})\\
    U\qty(\Lambda_1)\qty(\mathbbm 1+\frac\im 2{\omega_2}_{\mu\nu}M^{\mu\nu})U^{-1}\qty(\Lambda_1) &=U(\mathbbm 1+\Lambda_1 \omega_2 \Lambda_1^{\textnormal T})\\
    \mathbbm 1+U\qty(\Lambda_1)\frac\im 2{\omega_2}_{\mu\nu}M^{\mu\nu}U^{-1}\qty(\Lambda_1) &=\mathbbm 1 +\frac\im 2\tensor{\qty(\Lambda_1 \omega_2 \Lambda_1^{\textnormal T})}{_\alpha_\beta}M^{\alpha\beta}\\
    {\omega_2}_{\mu\nu}U\qty(\Lambda_1)M^{\mu\nu}U^{-1}\qty(\Lambda_1) &=\tensor{{\Lambda_1}}{_\alpha^\mu}\tensor{{\omega_2}}{_\mu_\nu}\tensor{{\Lambda_1}}{_\beta^\nu}M^{\alpha\beta}\\
    U\qty(\Lambda_1)M^{\mu\nu}U^{-1}\qty(\Lambda_1) &=\tensor{{\Lambda_1}}{_\alpha^\mu}\tensor{{\Lambda_1}}{_\beta^\nu}M^{\alpha\beta}
\end{align}
this equation already is useful by itself, as it shows how the generators themselves transforms under a group 
action, but, as we're interested into the algebra of the group, we still need to expand around $\Lambda_2$,
\begin{align}
    U\qty(\mathbbm 1+\omega_1)M^{\mu\nu}U^{-1}\qty(\mathbbm 1+\omega_1) &=\qty(\tensor{g}{_\alpha^\mu}+\tensor{{\omega_1}}{_\alpha^\mu})\qty(\tensor{g}{_\beta^\nu}+\tensor{{\omega_1}}{_\beta^\nu})M^{\alpha\beta}\\
    \qty(\mathbbm 1 +\frac\im 2{\omega_1}_{\alpha\beta}M^{\alpha\beta})M^{\mu\nu}\qty(\mathbbm 1 -\frac\im 2{\omega_1}_{\alpha\beta}M^{\alpha\beta})&=M^{\mu\nu}+M^{\mu\beta}\tensor{{\omega_1}}{_\beta^\nu}+M^{\alpha\nu}\tensor{{\omega_1}}{_\alpha^\mu}\\
    M^{\mu\nu}+\frac\im 2{\omega_1}_{\alpha\beta}M^{\alpha\beta}M^{\mu\nu} -\frac\im 2{\omega_1}_{\alpha\beta}M^{\mu\nu}M^{\alpha\beta}&=M^{\mu\nu}+M^{\mu\beta}\tensor{{\omega_1}}{_\beta^\nu}+M^{\alpha\nu}\tensor{{\omega_1}}{_\alpha^\mu}\\
    \frac\im 2{\omega_1}_{\alpha\beta}\comm{M^{\alpha\beta}}{M^{\mu\nu}}&={\omega_1}_{\alpha\beta}\qty(-\tensor{g}{^\alpha^\nu}M^{\mu\beta}+M^{\alpha\nu}g^{\beta\mu})\\
    \frac\im 2{\omega_1}_{\alpha\beta}\comm{M^{\alpha\beta}}{M^{\mu\nu}}&=\frac12{\omega_1}_{\alpha\beta}\qty(-\tensor{g}{^\alpha^\nu}M^{\mu\beta}+M^{\alpha\nu}g^{\beta\mu}+\tensor{g}{^\beta^\nu}M^{\mu\alpha}-M^{\beta\nu}g^{\alpha\mu})
\end{align}
thus,
\begin{align}
    \im \comm{M^{\alpha\beta}}{M^{\mu\nu}}&=-\tensor{g}{^\alpha^\nu}M^{\mu\beta}+M^{\alpha\nu}g^{\beta\mu}+\tensor{g}{^\beta^\nu}M^{\mu\alpha}-M^{\beta\nu}g^{\alpha\mu}
\end{align}

\probitem{}